\documentclass[12pt,a4paper]{article}
\usepackage[utf8]{inputenc}
\usepackage[T1]{fontenc}
\usepackage{amsmath}
\usepackage{amssymb}
\usepackage{graphicx}
\usepackage[french]{babel}
\usepackage[a4paper, margin=2.5cm]{geometry}
\usepackage[european]{circuitikz} % Pour dessiner des circuits
\usetikzlibrary{babel} % Pour la compatibilité avec babel-french
\usepackage{comment}

\title{Solution TP3 : Modèles de Thévenin et Norton}
\author{Prof: CERO}
\date{\today}

\begin{document}

\maketitle

\section*{Introduction}
Ce document présente la résolution détaillée des exercices du TP3 concernant les modèles 
de Thévenin et Norton, ainsi que le théorème de superposition. Les explications sont conçues 
pour être claires et didactiques. Le nœud 'A' (+) se trouve à la jonction des résistances de 
1.8k$\Omega$ et 1k$\Omega$, tandis que le nœud 'B' (-) est connecté à la masse.

\begin{center}
    \begin{circuitikz}
        \draw (0,0) node[ground]{} to[V, v=$15V$] (0,3) to[R, l=$3.9k\Omega$] (2,3);
        \draw (2,0) node[ground]{} to[V, v=$5V$] (2,1) to[R, l=$3.9k\Omega$] (2,3);
        \draw (4,0) node[ground]{} to[V, v=$15V$] (4,1) to[R, l=$1k\Omega$] (4,3);
        \draw (2,3) node[circ, label=above:$V_C$]{} to[R, l=$1.8k\Omega$] (4,3) node[circ, label=above:A]{};
        
        % Lignes pointillées pour délimiter le dipôle D
        \draw[dashed] (4, 3) -- (5, 3);
        %\node at (5.5, 3.7) {D};
        \node at (5.5, 3) {(+)};
        \node at (5.5, 0) {(-)};
        \draw (6,0) node[ground, label=right:B]{};
    \end{circuitikz}
\end{center}

\section{Préparation}

\subsection*{Rappel des modèles de Thévenin et Norton}
\begin{itemize}
    \item \textbf{Modèle de Thévenin :} Un circuit linéaire complexe peut être remplacé par une source de tension idéale $E_{Th}$ en série avec une résistance $R_{Th}$.
    \item \textbf{Modèle de Norton :} Le même circuit peut être remplacé par une source de courant idéale $I_N$ en parallèle avec une résistance $R_{N}$.
\end{itemize}
La résistance de Thévenin est égale à la résistance de Norton ($R_{Th} = R_N$). La relation entre les modèles est donnée par la loi d'Ohm : $E_{Th} = R_{Th} \cdot I_N$.

\subsection*{a) Différence de potentiel $V_{AB}$ par superposition}
Pour trouver la tension $V_{AB}$ (qui est la tension en circuit ouvert, donc $E_{Th}$), nous utilisons le théorème de 
superposition. Nous calculons la contribution de chaque source de tension indépendamment, en remplaçant les autres par un 
court-circuit.

Le point B est notre référence (masse, 0V). La tension $V_{AB}$ est donc égale à $V_A$.

Calculons $V_A$ en utilisant la superposition.

\paragraph{1. Contribution de la source de 15V (gauche), $V_1$ :}
Les autres sources (5V et 15V droite) sont court-circuitées. Le circuit devient :

\begin{center}
	\begin{circuitikz}
		\draw (0,0) node[ground]{} to[V, v=$15V$] (0,3) to[R, l=$3.9k\Omega$] (3,3) node[circ, label=above:$C$]{};
		\draw (3,3) to[R, l=$3.9k\Omega$] (3,0) node[ground]{};
		\draw (3,3) to[R, l=$1.8k\Omega$] (6,3) node[circ, label=above:$A$]{};
		\draw (6,3) to[R, l=$1k\Omega$] (6,0) node[ground]{};
	\end{circuitikz}
\end{center}

Pour résoudre ce circuit, nous utilisons l'analyse nodale pour les nœuds A et C.
Les équations de la Loi des Courants de Kirchhoff (LCK) sont :
\begin{itemize}
    \item \textbf{Nœud C :} $\frac{V_C - 15V}{3.9k\Omega} + \frac{V_C}{3.9k\Omega} + \frac{V_C - V_A}{1.8k\Omega} = 0$
    \item \textbf{Nœud A :} $\frac{V_A - V_C}{1.8k\Omega} + \frac{V_A}{1k\Omega} = 0$
\end{itemize}
De l'équation du nœud A, nous pouvons isoler $V_C$ :
$$ V_A \left( \frac{1}{1.8k} + \frac{1}{1k} \right) = \frac{V_C}{1.8k} \implies V_C = V_A \left( 1 + \frac{1.8k}{1k} \right) = 2.8 \cdot V_A $$
En substituant $V_C$ dans l'équation du nœud C (et en multipliant tout par $3.9k \cdot 1.8k$ pour simplifier) :
$$ 1.8(2.8 V_A - 15) + 1.8(2.8 V_A) + 3.9(2.8 V_A - V_A) = 0 $$
$$ 5.04 V_A - 27 + 5.04 V_A + 3.9(1.8 V_A) = 0 $$
$$ 10.08 V_A - 27 + 7.02 V_A = 0 $$
$$ 17.1 V_A = 27 $$
$$ V_{A1} = \frac{27}{17.1} \approx \mathbf{1.579 \, V} $$

\paragraph{2. Contribution de la source de 5V (centre), $V_2$ :}
Les deux sources de 15V sont court-circuitées.
\begin{center}
\begin{circuitikz}
    \draw (0,0) node[ground]{} to[R, l=$3.9k\Omega$] (0,3) to (3,3) node[circ, label=above:$C$]{};
    \draw (3,0) node[ground]{} to[V, v=$5V$] (3,1) to[R, l=$3.9k\Omega$] (3,3);
    \draw (3,3) to[R, l=$1.8k\Omega$] (6,3) node[circ, label=above:$A$]{};
    \draw (6,3) to[R, l=$1k\Omega$] (6,0) node[ground]{};
\end{circuitikz}
\end{center}

Pour analyser ce circuit, nous transformons d'abord la source de 5V et sa résistance en série de 3.9k$\Omega$ en son équivalent de Norton :
$$ I_{N,5V} = \frac{5V}{3.9k\Omega} \approx 1.282 \, mA \quad \text{avec une résistance en parallèle de } 3.9k\Omega $$
Cette source de courant est connectée au nœud C. Au nœud C, nous avons aussi la résistance de 3.9k$\Omega$ de la branche de gauche. La résistance équivalente de ces deux en parallèle est $R_{eq,C} = \frac{3.9k}{2} = 1.95k\Omega$.

Le circuit se simplifie en une source de courant $I_{N,5V}$ qui se divise entre $R_{eq,C}$ et la résistance en série de la branche du nœud A ($1.8k\Omega + 1k\Omega = 2.8k\Omega$).

Nous calculons le courant $I_A$ qui passe par la branche du nœud A en utilisant le diviseur de courant :
$$ I_A = I_{N,5V} \cdot \frac{R_{eq,C}}{R_{eq,C} + (1.8k\Omega + 1k\Omega)} $$
$$ I_A = 1.282 \, mA \cdot \frac{1.95k\Omega}{1.95k\Omega + 2.8k\Omega} = 1.282 \, mA \cdot \frac{1950}{4750} \approx 0.526 \, mA $$
Finalement, la tension au nœud A est ce courant multiplié par la résistance de 1k$\Omega$ :
$$ V_{A2} = I_A \cdot 1k\Omega = 0.526 \, mA \cdot 1k\Omega = \mathbf{0.526 \, V} $$

\paragraph{3. Contribution de la source de 15V (droite), $V_3$ :}
Les sources de 15V (gauche) et 5V sont court-circuitées.
\begin{center}
\begin{circuitikz}
    \draw (0,0) node[ground]{} to[R, l=$3.9k\Omega$] (0,3) to (3,3) node[circ, label=above:$C$]{};
    \draw (3,0) node[ground]{} to[R, l=$3.9k\Omega$] (3,3);
    \draw (3,3) to[R, l=$1.8k\Omega$] (6,3) node[circ, label=above:$A$]{};
    \draw (6,0) node[ground]{} to[V, v=$15V$] (6,1) to[R, l=$1k\Omega$] (6,3);
\end{circuitikz}
\end{center}

En court-circuitant les sources de gauche et du centre, leurs résistances de 3.9k$\Omega$ se retrouvent en parallèle entre le nœud C et la masse. Leur équivalent est $R_{eq,C} = \frac{3.9k\Omega}{2} = 1.95k\Omega$.

La manière la plus simple de le résoudre est de transformer la source de 15V et sa résistance de 1k$\Omega$ en son équivalent de Norton :
$$ I_{N,15V} = \frac{15V}{1k\Omega} = 15 \, mA \quad \text{avec une résistance en parallèle de } 1k\Omega $$
Cette source de courant est connectée au nœud A. La tension en A sera le produit de ce courant par la résistance totale équivalente vue depuis A. La résistance vue depuis A est celle de 1k$\Omega$ en parallèle avec le chemin qui va vers C, qui est $1.8k\Omega + R_{eq,C} = 1.8k\Omega + 1.95k\Omega = 3.75k\Omega$.

La résistance totale équivalente au nœud A est :
$$ R_{eq,A} = \frac{1k\Omega \cdot 3.75k\Omega}{1k\Omega + 3.75k\Omega} = \frac{3750}{4.75} \approx 789.5 \, \Omega $$
La tension en A est :
$$ V_{A3} = I_{N,15V} \cdot R_{eq,A} = 15 \, mA \cdot 789.5 \, \Omega = \mathbf{11.842 \, V} $$

\paragraph{Tension totale :}
La tension totale au nœud A, qui est la tension de Thévenin $E_{Th}$, est la somme des trois contributions calculées :
$$ E_{Th} = V_{A1} + V_{A2} + V_{A3} $$
$$ E_{Th} = 1.579 \, V + 0.526 \, V + 11.842 \, V = \mathbf{13.947 \, V} $$
\textit{Note : Cette valeur est le résultat correct pour le circuit analysé. L'écart avec la valeur de 10.8V de l'énoncé original persiste, ce qui confirme qu'il y a une erreur dans les données du problème ou dans le schéma original que l'on tentait de représenter.}

\subsection*{b) Détermination de $E_{Th}$, $I_N$ et $R_{Th}$}
\paragraph{$E_{Th}$ (Tension de Thévenin) :}
La tension de Thévenin est la tension à vide entre A et B, que nous venons de calculer.
$$ E_{Th} = V_{AB} = \mathbf{13.947 \, V} $$

\paragraph{$R_{Th}$ (Résistance de Thévenin) :}
Pour trouver $R_{Th}$, on 'éteint' toutes les sources de tension (on les remplace par un 
court-circuit) et on calcule la résistance équivalente vue depuis les bornes A et B.
\begin{center}
    \begin{circuitikz}
        \draw (0,0) node[ground]{} to[R, l=$3.9k\Omega$] (0,3) to (2,3) node[circ, label=above:$C$]{};
        \draw (2,0) node[ground]{} to[R, l=$3.9k\Omega$] (2,3);
        \draw (2,3) to[R, l=$1.8k\Omega$] (5,3) node[circ, label=right:$A$]{};
        \draw (5,3) to[R, l=$1k\Omega$] (5,0) node[ground, label=right:B]{};
        \node at (6,3) {}; % Point de mesure
    \end{circuitikz}
\end{center}
Les deux résistances de 3.9k$\Omega$ sont en parallèle, leur équivalent est $R_{eq,C} = 1.95k\Omega$. Cette résistance est en série avec celle de 1.8k$\Omega$. L'ensemble est en parallèle avec la résistance de 1k$\Omega$.
$$ R_{Th} = (1.95k\Omega + 1.8k\Omega) \parallel 1k\Omega = 3.75k\Omega \parallel 1k\Omega $$
$$ R_{Th} = \frac{3.75k\Omega \cdot 1k\Omega}{3.75k\Omega + 1k\Omega} = \frac{3750}{4.75} \approx \mathbf{789.5 \, \Omega} $$
\textit{Note : Cette valeur coïncide avec la valeur de référence de 790$\Omega$ de l'énoncé, ce qui confirme que la topologie du circuit est correcte. Cependant, la valeur de $E_{Th}$ que nous avons calculée (13.95V) ne coïncide pas avec celle de l'énoncé (10.8V).}

\paragraph{$I_N$ (Courant de Norton) :}
Le courant de Norton est le courant qui circule dans un court-circuit placé entre A et B. On peut le calculer avec la loi d'Ohm à partir de $E_{Th}$ et $R_{Th}$.
$$ I_N = \frac{E_{Th}}{R_{Th}} = \frac{13.947 \, V}{789.5 \, \Omega} \approx 0.01766 \, A = \mathbf{17.67 \, mA} $$
\textit{Note : Cette valeur est proche de la référence de l'énoncé (13.7mA), bien qu'elle ne soit pas identique en raison de l'écart sur $E_{Th}$.}

\subsection*{c) Tension avec une charge de 790 $\Omega$}
Nous branchons une résistance de charge $R_L = 790 \, \Omega$ entre A et B. En utilisant le modèle de Thévenin, le circuit devient un simple diviseur de tension.
\begin{center}
    \begin{circuitikz}
        \draw (0,0) node[ground]{} to[V, v=$E_{Th}$] (0,2) to[R, l=$R_{Th}$] (2,2) to[R, l={$R_L=790\Omega$}] (2,0) node[ground]{};
        \node at (2,2.3) {$V_L$};
    \end{circuitikz}
\end{center}
La tension $V_L$ aux bornes de la charge est :
$$ V_L = E_{Th} \cdot \frac{R_L}{R_{Th} + R_L} $$
En utilisant nos valeurs calculées :
$$ V_L = 13.947 \cdot \frac{790}{789.5 + 790} = 13.947 \cdot \frac{790}{1579.5} \approx \mathbf{6.97 \, V} $$
Si nous utilisions les valeurs (incorrectes) de l'énoncé ($E_{Th}=10.8V, R_{Th}=790\Omega$) :
$$ V_L = 10.8 \cdot \frac{790}{790 + 790} = 10.8 \cdot \frac{1}{2} = 5.4 \, V $$
Ceci montre que lorsque $R_L = R_{Th}$, la tension aux bornes de la charge est la moitié de la tension de Thévenin.

\end{document}
